\documentclass[11pt]{article}
\usepackage[margin=1in]{geometry}
\usepackage{amsmath, amssymb}
\usepackage{booktabs}
\usepackage{graphicx}
\usepackage{hyperref}
\usepackage{caption}
\usepackage{subcaption}
\usepackage{multirow}
\usepackage{array}
\usepackage{enumitem}
\usepackage{xcolor}
\usepackage{listings}

\lstset{
    language=Python,
    basicstyle=\ttfamily\small,
    keywordstyle=\color{blue},
    commentstyle=\color{green!40!black},
    stringstyle=\color{red},
    showstringspaces=false,
    breaklines=true,
    frame=single,
    numbers=left,
    numbersep=5pt,
}

\title{Detecting Adversarial Perturbations in AI-Generated Code:\\A Study of Static and Embedding-Based Approaches}

\author{Anonymous Author(s)}

\date{}

\begin{document}
\maketitle

\begin{abstract}
% ... (keep existing abstract)
\end{abstract}

\section{Introduction}
\label{sec:intro}
% ... (keep existing + add contributions explicitly)

\textbf{Research Questions:}
\begin{itemize}[itemsep=0pt, topsep=0pt]
    \item \textbf{RQ1:} What perturbation strategies can adversarially modify AI-generated code while preserving syntactic plausibility?
    \item \textbf{RQ2:} How effectively do current detection approaches (static AST features and learned embeddings) identify these perturbations?
    \item \textbf{RQ3:} Is there a fundamental limit to what static and embedding-based methods can detect?
\end{itemize}

\textbf{Contributions:}
\begin{itemize}[itemsep=0pt, topsep=0pt]
    \item A systematic taxonomy of five perturbation strategies for AI-generated code
    \item A curated dataset of 550 labeled Python functions
    \item A comparative evaluation of AST-based and embedding-based detectors
    \item The discovery of a three-tier detectability hierarchy
    \item Statistical validation using McNemar's test
\end{itemize}

% ... (rest of sections stay same, with these additions)

\subsection{Statistical Analysis}
\label{sec:statistical}

To determine whether the performance difference between Track A (AST features) and Track B (CodeBERT embeddings) is statistically significant, we apply McNemar's test with continuity correction. The contingency table is:

\begin{table}[h]
\centering
\begin{tabular}{@{}lcc@{}}
\toprule
& \textbf{AST Correct} & \textbf{AST Incorrect} \\
\midrule
\textbf{Embedding Correct} & 475 & 9 \\
\textbf{Embedding Incorrect} & 23 & 43 \\
\bottomrule
\end{tabular}
\caption{Contingency table for McNemar's test comparing AST and CodeBERT detectors.}
\label{tab:mcnemar}
\end{table}

The test yields $\chi^2 = 9.00$ with $p = 0.0201$, confirming that the AST-based detector significantly outperforms the embedding-based detector at $\alpha = 0.05$.

\end{document}
